% res_mlp_diagram.tex
% Diagram reziduálneho MLP modelu s vysvetlením LayerNorm a Dropout
% Použitie: \resizebox{0.80\textwidth}{!}{% res_mlp_diagram.tex
% Diagram reziduálneho MLP modelu s vysvetlením LayerNorm a Dropout
% Použitie: \resizebox{0.80\textwidth}{!}{% res_mlp_diagram.tex
% Diagram reziduálneho MLP modelu s vysvetlením LayerNorm a Dropout
% Použitie: \resizebox{0.80\textwidth}{!}{% res_mlp_diagram.tex
% Diagram reziduálneho MLP modelu s vysvetlením LayerNorm a Dropout
% Použitie: \resizebox{0.80\textwidth}{!}{\input{figures/tikz/res_mlp_diagram}}
\begin{tikzpicture}[
    font=\small,
    box/.style={
        draw, rounded corners=5pt, align=center,
        minimum width=3.4cm, minimum height=0.82cm,
        inner sep=4pt
    },
    inputbox/.style={box, fill=gray!10, draw=gray!50},
    projbox/.style={box, fill=blue!8, draw=blue!40},
    normbox/.style={box, fill=yellow!10, draw=yellow!60!black},
    dropbox/.style={box, fill=red!6, draw=red!30},
    outbox/.style={box, fill=green!8, draw=green!50},
    labelbox/.style={font=\scriptsize\itshape, text=gray!65, align=left,
                     text width=4.2cm},
    arrow/.style={-{Latex[length=2.5mm]}, thick},
    resarrow/.style={-{Latex[length=2.5mm]}, thick, draw=green!55!black}
]

\node[inputbox] (phi) at (0, 0)
    {\(\phi_{ij} \in \mathbb{R}^{10}\)\\[-2pt]
     \scriptsize vektor príznakov hrany};

\node[projbox] (proj) at (0, -1.55)
    {Vstupná projekcia\\[-2pt]
     \scriptsize + aktivácia GELU};

\node[normbox] (ln1) at (0, -3.10)
    {Vrstvová normalizácia\\[-2pt]
     \scriptsize (LayerNorm)};

\node[projbox] (fc1) at (0, -4.50)
    {Lineárna transformácia\\[-2pt]
     \scriptsize + aktivácia GELU};

\node[dropbox] (drop) at (0, -5.85)
    {Náhodné vypúšťanie\\[-2pt]
     \scriptsize (Dropout)};

\node[projbox] (fc2) at (0, -7.20)
    {Lineárna transformácia};

\node[draw=green!55!black, fill=green!5, circle,
      minimum size=0.65cm, font=\bfseries\small] (plus) at (0, -8.55) {+};

\draw[draw=green!40, dashed, rounded corners=6pt]
    (-2.3, -2.55) rectangle (2.3, -7.85);
\node[font=\scriptsize\itshape, text=green!60!black]
    at (0, -2.38) {reziduálny blok (opakuje sa \(D\)-krát)};

\draw[resarrow]
    (proj.east) -- ++(1.4, 0) -- ++(0, -7.0) -- (plus.east);
\node[font=\scriptsize, text=green!60!black, rotate=90, anchor=south]
    at (3.0, -5.2) {reziduálna skratka};

\node[outbox] (out) at (0, -9.95)
    {Záverečná lineárna vrstva\\[-2pt]
     \scriptsize logit \(s_{ij} \in \mathbb{R}\)};

\draw[arrow] (phi)  -- (proj);
\draw[arrow] (proj) -- (ln1);
\draw[arrow] (ln1)  -- (fc1);
\draw[arrow] (fc1)  -- (drop);
\draw[arrow] (drop) -- (fc2);
\draw[arrow] (fc2)  -- (plus);
\draw[arrow] (plus) -- (out);

\node[labelbox, anchor=north west] at (3.5, -3.25)
    {\textbf{LayerNorm:} normalizuje aktivácie
     v rámci jedného vzorku — stabilizuje tréning
     a urýchľuje konvergenciu.};

\node[labelbox, anchor=north west] at (3.5, -5.85)
    {\textbf{Dropout:} náhodne deaktivuje
     časť neurónov — zabraňuje preučeniu.};

\node[labelbox, anchor=north west] at (3.5, -8.55)
    {\textbf{Reziduálna skratka:}
     gradient prúdi priamo, bez
     útlmu cez hlboké bloky.};

\end{tikzpicture}
}
\begin{tikzpicture}[
    font=\small,
    box/.style={
        draw, rounded corners=5pt, align=center,
        minimum width=3.4cm, minimum height=0.82cm,
        inner sep=4pt
    },
    inputbox/.style={box, fill=gray!10, draw=gray!50},
    projbox/.style={box, fill=blue!8, draw=blue!40},
    normbox/.style={box, fill=yellow!10, draw=yellow!60!black},
    dropbox/.style={box, fill=red!6, draw=red!30},
    outbox/.style={box, fill=green!8, draw=green!50},
    labelbox/.style={font=\scriptsize\itshape, text=gray!65, align=left,
                     text width=4.2cm},
    arrow/.style={-{Latex[length=2.5mm]}, thick},
    resarrow/.style={-{Latex[length=2.5mm]}, thick, draw=green!55!black}
]

\node[inputbox] (phi) at (0, 0)
    {\(\phi_{ij} \in \mathbb{R}^{10}\)\\[-2pt]
     \scriptsize vektor príznakov hrany};

\node[projbox] (proj) at (0, -1.55)
    {Vstupná projekcia\\[-2pt]
     \scriptsize + aktivácia GELU};

\node[normbox] (ln1) at (0, -3.10)
    {Vrstvová normalizácia\\[-2pt]
     \scriptsize (LayerNorm)};

\node[projbox] (fc1) at (0, -4.50)
    {Lineárna transformácia\\[-2pt]
     \scriptsize + aktivácia GELU};

\node[dropbox] (drop) at (0, -5.85)
    {Náhodné vypúšťanie\\[-2pt]
     \scriptsize (Dropout)};

\node[projbox] (fc2) at (0, -7.20)
    {Lineárna transformácia};

\node[draw=green!55!black, fill=green!5, circle,
      minimum size=0.65cm, font=\bfseries\small] (plus) at (0, -8.55) {+};

\draw[draw=green!40, dashed, rounded corners=6pt]
    (-2.3, -2.55) rectangle (2.3, -7.85);
\node[font=\scriptsize\itshape, text=green!60!black]
    at (0, -2.38) {reziduálny blok (opakuje sa \(D\)-krát)};

\draw[resarrow]
    (proj.east) -- ++(1.4, 0) -- ++(0, -7.0) -- (plus.east);
\node[font=\scriptsize, text=green!60!black, rotate=90, anchor=south]
    at (3.0, -5.2) {reziduálna skratka};

\node[outbox] (out) at (0, -9.95)
    {Záverečná lineárna vrstva\\[-2pt]
     \scriptsize logit \(s_{ij} \in \mathbb{R}\)};

\draw[arrow] (phi)  -- (proj);
\draw[arrow] (proj) -- (ln1);
\draw[arrow] (ln1)  -- (fc1);
\draw[arrow] (fc1)  -- (drop);
\draw[arrow] (drop) -- (fc2);
\draw[arrow] (fc2)  -- (plus);
\draw[arrow] (plus) -- (out);

\node[labelbox, anchor=north west] at (3.5, -3.25)
    {\textbf{LayerNorm:} normalizuje aktivácie
     v rámci jedného vzorku — stabilizuje tréning
     a urýchľuje konvergenciu.};

\node[labelbox, anchor=north west] at (3.5, -5.85)
    {\textbf{Dropout:} náhodne deaktivuje
     časť neurónov — zabraňuje preučeniu.};

\node[labelbox, anchor=north west] at (3.5, -8.55)
    {\textbf{Reziduálna skratka:}
     gradient prúdi priamo, bez
     útlmu cez hlboké bloky.};

\end{tikzpicture}
}
\begin{tikzpicture}[
    font=\small,
    box/.style={
        draw, rounded corners=5pt, align=center,
        minimum width=3.4cm, minimum height=0.82cm,
        inner sep=4pt
    },
    inputbox/.style={box, fill=gray!10, draw=gray!50},
    projbox/.style={box, fill=blue!8, draw=blue!40},
    normbox/.style={box, fill=yellow!10, draw=yellow!60!black},
    dropbox/.style={box, fill=red!6, draw=red!30},
    outbox/.style={box, fill=green!8, draw=green!50},
    labelbox/.style={font=\scriptsize\itshape, text=gray!65, align=left,
                     text width=4.2cm},
    arrow/.style={-{Latex[length=2.5mm]}, thick},
    resarrow/.style={-{Latex[length=2.5mm]}, thick, draw=green!55!black}
]

\node[inputbox] (phi) at (0, 0)
    {\(\phi_{ij} \in \mathbb{R}^{10}\)\\[-2pt]
     \scriptsize vektor príznakov hrany};

\node[projbox] (proj) at (0, -1.55)
    {Vstupná projekcia\\[-2pt]
     \scriptsize + aktivácia GELU};

\node[normbox] (ln1) at (0, -3.10)
    {Vrstvová normalizácia\\[-2pt]
     \scriptsize (LayerNorm)};

\node[projbox] (fc1) at (0, -4.50)
    {Lineárna transformácia\\[-2pt]
     \scriptsize + aktivácia GELU};

\node[dropbox] (drop) at (0, -5.85)
    {Náhodné vypúšťanie\\[-2pt]
     \scriptsize (Dropout)};

\node[projbox] (fc2) at (0, -7.20)
    {Lineárna transformácia};

\node[draw=green!55!black, fill=green!5, circle,
      minimum size=0.65cm, font=\bfseries\small] (plus) at (0, -8.55) {+};

\draw[draw=green!40, dashed, rounded corners=6pt]
    (-2.3, -2.55) rectangle (2.3, -7.85);
\node[font=\scriptsize\itshape, text=green!60!black]
    at (0, -2.38) {reziduálny blok (opakuje sa \(D\)-krát)};

\draw[resarrow]
    (proj.east) -- ++(1.4, 0) -- ++(0, -7.0) -- (plus.east);
\node[font=\scriptsize, text=green!60!black, rotate=90, anchor=south]
    at (3.0, -5.2) {reziduálna skratka};

\node[outbox] (out) at (0, -9.95)
    {Záverečná lineárna vrstva\\[-2pt]
     \scriptsize logit \(s_{ij} \in \mathbb{R}\)};

\draw[arrow] (phi)  -- (proj);
\draw[arrow] (proj) -- (ln1);
\draw[arrow] (ln1)  -- (fc1);
\draw[arrow] (fc1)  -- (drop);
\draw[arrow] (drop) -- (fc2);
\draw[arrow] (fc2)  -- (plus);
\draw[arrow] (plus) -- (out);

\node[labelbox, anchor=north west] at (3.5, -3.25)
    {\textbf{LayerNorm:} normalizuje aktivácie
     v rámci jedného vzorku — stabilizuje tréning
     a urýchľuje konvergenciu.};

\node[labelbox, anchor=north west] at (3.5, -5.85)
    {\textbf{Dropout:} náhodne deaktivuje
     časť neurónov — zabraňuje preučeniu.};

\node[labelbox, anchor=north west] at (3.5, -8.55)
    {\textbf{Reziduálna skratka:}
     gradient prúdi priamo, bez
     útlmu cez hlboké bloky.};

\end{tikzpicture}
}
\begin{tikzpicture}[
    font=\small,
    box/.style={
        draw, rounded corners=5pt, align=center,
        minimum width=3.4cm, minimum height=0.82cm,
        inner sep=4pt
    },
    inputbox/.style={box, fill=gray!10, draw=gray!50},
    projbox/.style={box, fill=blue!8, draw=blue!40},
    normbox/.style={box, fill=yellow!10, draw=yellow!60!black},
    dropbox/.style={box, fill=red!6, draw=red!30},
    outbox/.style={box, fill=green!8, draw=green!50},
    labelbox/.style={font=\scriptsize\itshape, text=gray!65, align=left,
                     text width=4.2cm},
    arrow/.style={-{Latex[length=2.5mm]}, thick},
    resarrow/.style={-{Latex[length=2.5mm]}, thick, draw=green!55!black}
]

\node[inputbox] (phi) at (0, 0)
    {\(\phi_{ij} \in \mathbb{R}^{10}\)\\[-2pt]
     \scriptsize vektor príznakov hrany};

\node[projbox] (proj) at (0, -1.55)
    {Vstupná projekcia\\[-2pt]
     \scriptsize + aktivácia GELU};

\node[normbox] (ln1) at (0, -3.10)
    {Vrstvová normalizácia\\[-2pt]
     \scriptsize (LayerNorm)};

\node[projbox] (fc1) at (0, -4.50)
    {Lineárna transformácia\\[-2pt]
     \scriptsize + aktivácia GELU};

\node[dropbox] (drop) at (0, -5.85)
    {Náhodné vypúšťanie\\[-2pt]
     \scriptsize (Dropout)};

\node[projbox] (fc2) at (0, -7.20)
    {Lineárna transformácia};

\node[draw=green!55!black, fill=green!5, circle,
      minimum size=0.65cm, font=\bfseries\small] (plus) at (0, -8.55) {+};

\draw[draw=green!40, dashed, rounded corners=6pt]
    (-2.3, -2.55) rectangle (2.3, -7.85);
\node[font=\scriptsize\itshape, text=green!60!black]
    at (0, -2.38) {reziduálny blok (opakuje sa \(D\)-krát)};

\draw[resarrow]
    (proj.east) -- ++(1.4, 0) -- ++(0, -7.0) -- (plus.east);
\node[font=\scriptsize, text=green!60!black, rotate=90, anchor=south]
    at (3.0, -5.2) {reziduálna skratka};

\node[outbox] (out) at (0, -9.95)
    {Záverečná lineárna vrstva\\[-2pt]
     \scriptsize logit \(s_{ij} \in \mathbb{R}\)};

\draw[arrow] (phi)  -- (proj);
\draw[arrow] (proj) -- (ln1);
\draw[arrow] (ln1)  -- (fc1);
\draw[arrow] (fc1)  -- (drop);
\draw[arrow] (drop) -- (fc2);
\draw[arrow] (fc2)  -- (plus);
\draw[arrow] (plus) -- (out);

\node[labelbox, anchor=north west] at (3.5, -3.25)
    {\textbf{LayerNorm:} normalizuje aktivácie
     v rámci jedného vzorku — stabilizuje tréning
     a urýchľuje konvergenciu.};

\node[labelbox, anchor=north west] at (3.5, -5.85)
    {\textbf{Dropout:} náhodne deaktivuje
     časť neurónov — zabraňuje preučeniu.};

\node[labelbox, anchor=north west] at (3.5, -8.55)
    {\textbf{Reziduálna skratka:}
     gradient prúdi priamo, bez
     útlmu cez hlboké bloky.};

\end{tikzpicture}
